\section{Introduction}

\lettrine[lines=3,lhang=0.08,loversize=0.15,findent=0.5em,nindent=0em]{T}{he} operational security of symmetric cryptography depends on secret keys. When an encryption key is stored in a single database, device, or password vault, the entire protection mechanism depends on the continued availability and integrity of that one location. The failure modes are immediate: key loss makes ciphertext permanently unreadable, while key compromise allows an attacker to decrypt all protected data. For this reason, practical key management must address both confidentiality and recoverability.

Threshold secret sharing offers a direct way to remove this single point of failure. Shamir's scheme divides a secret into multiple shares such that any threshold number of shares can reconstruct the secret, while any smaller set reveals no information about it \cite{shamir1979}. The construction reduces the access-control problem to polynomial interpolation over a finite field. However, standard reconstruction assumes that the shares presented by participants are correct. In realistic settings this assumption is weak: shares may be unavailable, corrupted by storage faults, or intentionally modified by an adversary.

The close algebraic relationship between Shamir sharing and Reed-Solomon codes makes it natural to strengthen reconstruction with error correction \cite{reedsolomon1960}. Both mechanisms evaluate a low-degree polynomial at distinct field points. As a result, robust recovery can be formulated as recovering the original polynomial from noisy point evaluations, which is precisely the setting addressed by Reed-Solomon decoding algorithms. In particular, the Berlekamp-Welch perspective allows reconstruction to remain correct even when a bounded number of submitted shares are wrong \cite{fedorenko2005}.

This paper studies a practical key recovery prototype rather than a new cryptographic primitive. The protected secret is an AES-256 key \cite{fips197}, used with the AES-GCM authenticated encryption mode \cite{sp80038d}. The contribution lies in implementing a modular end-to-end system and experimentally validating the threshold and error-correction behavior of the resulting recovery process.

The main contributions of this work are as follows:
\begin{enumerate}
\item A modular Python implementation that integrates AES-GCM encryption, Shamir share generation, and Reed-Solomon-style Berlekamp-Welch reconstruction into a single key recovery workflow.
\item An experimental evaluation over fifteen recovery scenarios that cover normal recovery, minimum-threshold recovery, insufficient-share failure, and recovery with corrupted shares.
\item A grounded analysis showing that robust decoding improves availability and fault tolerance without reducing the secrecy threshold of the original sharing scheme.
\end{enumerate}

The remainder of the paper is organized as follows. Section II reviews the necessary theoretical background on AES-GCM, finite-field arithmetic, Shamir sharing, Reed-Solomon codes, and Berlekamp-Welch decoding. Section III presents the proposed recovery framework. Section IV describes the implementation. Section V defines the experiments, while Section VI reports the results and discusses their implications. Section VII analyzes the security model and limitations. Section VIII concludes the paper.
